\chapter{Certified polynomial inequalities}
\label{chap:certificates}

The odd--odd and mixed sectors rely on explicit two-variable inequalities
involving polynomials and $\log$. They are reduced to positivity of
polynomials on boxes, certified by nonnegative Bernstein coefficients after
subdivision. This chapter is independent of the rest of the analysis, so it
can be formalized in parallel. The natural Lean design is a verified Bernstein
checker (a function on lists of rational coefficients, proved sound once) run
by \texttt{decide} or \texttt{native\_decide}.

\section{General tools}

\begin{definition}[Bernstein coefficients]
\label{def:bernstein}
For a polynomial $\sum_{k\le n,\,l\le m}a_{kl}\xi^k\eta^l$ of bidegree $(n,m)$,
its Bernstein coefficients on $[0,1]^2$ are
\[
  b_{ij}=\sum_{k\le i,\,l\le j}a_{kl}\frac{\binom ik}{\binom nk}\frac{\binom jl}{\binom ml},
  \qquad 0\le i\le n,\ 0\le j\le m .
\]
On a box $[x_0,x_1]\times[y_0,y_1]$ one first applies the affine substitution
onto $[0,1]^2$.
\end{definition}

\begin{lemma}[Bernstein positivity criterion]
\label{lem:bernstein-positivity}
\uses{def:bernstein}
If all Bernstein coefficients of $P$ on a box are nonnegative, then $P\ge0$ on
that box. If moreover they are positive except for those with indices in a set
$Z$, then $P>0$ away from the corresponding faces and corners.
\end{lemma}

\begin{proof}
$P=\sum b_{ij}B^n_i(\xi)B^m_j(\eta)$ with nonnegative Bernstein basis
polynomials.
\end{proof}

\begin{lemma}[Rational bounds for the logarithm]
\label{lem:log-bounds}
For $0<r\le2$,
\[
  \frac{r-2}r\le\log(r/2)\le\frac{2(r-2)}{r+2}+\frac{2(r-2)^3}{3(r+2)^3},
\]
and for $0\le z\le\tfrac12$, with $T(z)=-\sum_{m=1}^{14}z^m/m$,
$T(z)-\frac{z^{15}}{15(1-z)}\le\log(1-z)\le T(z)$. Also
$\ell_-\le\log2\le\ell_+$ for $\ell_-=2\sum_{m=0}^{19}\frac{3^{-2m-1}}{2m+1}$
and $\ell_+=\ell_-+\frac{2\cdot3^{-41}}{41(1-3^{-2})}$.
\end{lemma}

\begin{proof}
Series expansions of $\log$ and $\operatorname{artanh}$ with explicit
remainders.
\end{proof}

\begin{lemma}[Affine dependence on the logarithm]
\label{lem:affine-in-log}
If $A+B\ell\ge0$ for $\ell=\ell_-$ and for $\ell=\ell_+$, then $A+B\ell\ge0$ for
every $\ell\in[\ell_-,\ell_+]$, regardless of the sign of $B$.
\end{lemma}

\begin{proof}
An affine function on an interval attains its minimum at an endpoint.
\end{proof}

\section{The certificates}

\begin{lemma}[Odd--odd residual]
\label{cert:dd-residual}
\uses{def:dd-ground-state}
$Pw_*+\tfrac52s^2w_*\ge0$ on the geometric range $0<s<2$,
$s(2-s)\le r\le2-s^2/2$.
\end{lemma}

\begin{proof}
\uses{lem:bernstein-positivity, lem:log-bounds, lem:affine-in-log, lem:dipole-identities}
Write $Pw_*+\tfrac52s^2w_*=R_0+R_1\log(r/2)$ with
$R_0=s\mathsf A(s,r)/(384r^3)$ and $R_1=s\mathsf B(s,r)/(1536r^2)$, where
$\mathsf A,\mathsf B$ are the integer polynomials in the manuscript's table
(the table was checked against a direct symbolic computation of
$Pw_*+\tfrac52s^2w_*$ at sample points). Using the two rational logarithm
bounds and the substitution $s=2x$, $r=2(1-x)[2x+(1-x)y]$, the resulting
polynomials are divisible by $(1-x)^4$; the quotients $Q_-$, $Q_+$ (bidegrees
$(9,6)$ and $(15,9)$) have nonnegative Bernstein coefficients on the six boxes
listed in the manuscript.
\end{proof}

\begin{lemma}[Odd-in-$b$ interior inequality]
\label{cert:odd-b}
\uses{def:odd-b-comparison}
For the comparison function of Definition~\ref{def:odd-b-comparison},
$\mathcal Pw_v\ge\tfrac1{100}h^2w_v$ on $(0,1]\times[0,1]$ in the $(h,k)$
variables.
\end{lemma}

\begin{proof}
\uses{lem:bernstein-positivity, lem:log-bounds, lem:affine-in-log}
After multiplying by $S^3R$ the inequality reads $A+B\log h\ge0$ with rational
polynomials $A,B$ of bidegrees $(10,8)$ and $(8,7)$. On $0\le h\le2^{-10}$ the
Bernstein coefficients of $-B$ and $A-6B$ are nonnegative. On the dyadic
strips $2^{-j-1}\le h\le2^{-j}$, $j=0,\dots,9$, use the logarithm bounds and
subdivision (658 nodes in total).
\end{proof}

\begin{lemma}[Odd-in-$a$ interior inequality]
\label{cert:odd-a}
\uses{def:odd-a-comparison}
For the comparison function of Definition~\ref{def:odd-a-comparison},
$\mathcal Pw_u\ge\tfrac1{200}h^2w_u$ on $Q$.
\end{lemma}

\begin{proof}
\uses{lem:bernstein-positivity}
Multiplying by $S^{16}R$ gives a rational polynomial $N(h,k)$ of bidegree
$(32,30)$. On the two triangles of the square, $t^{-15}N(tq,t)$ and
$t^{-15}N(t,tq)$ are polynomials of bidegrees $(17,32)$ and $(17,30)$ with
nonnegative Bernstein coefficients on two, respectively four, boxes.
\end{proof}

\begin{lemma}[Odd-in-$a$ Robin inequality]
\label{cert:robin}
\uses{def:odd-a-comparison}
$3h-2(1-h^2)\partial_hF_u(h,0)-(2-h^2)\partial_kF_u(h,0)-\tfrac h{250}\ge0$ for $h\in[0,1]$; the
left side is
$\tfrac1{1000}(673h^7-1953h^6-3213h^5+7317h^4+109h^3-4647h^2+1718h+1)-\tfrac h{250}$.
\end{lemma}

\begin{proof}
\uses{lem:bernstein-positivity}
One-variable Bernstein coefficients after dyadic subdivision of $[0,1]$ (nine
nodes, depth at most four).
\end{proof}

\begin{remark}[Digest note]
The manuscript refers to the scripts \texttt{verify\_odd\_b.py},
\texttt{verify\_odd\_a.py} and \texttt{verify\_robin.py}, which are not in this
repository. The polynomials $A,B$ (odd-in-$b$) and $N$ (odd-in-$a$) are only
described as ``generated by'' differentiation formulas. For formalization we
should either generate them inside Lean from the definitions of $F$, $\mathsf G$
and $U$ (most trustworthy), or import the coefficient lists and prove in Lean
that they agree with the defining expressions.
\end{remark}
